\title{Controlling Text-to-Image Diffusion by Orthogonal Finetuning}

\begin{abstract}
State-of-the-art text-to-image diffusion models excel at creating photorealistic imagery from textual descriptions. However, devising strategies to steer or control these versatile models for specialized downstream applications remains an open research issue. In response, we present Orthogonal Finetuning (OFT), a theoretically grounded method for finetuning text-to-image diffusion models to accommodate downstream requirements. In contrast to conventional approaches, OFT guarantees the preservation of hyperspherical energy, a metric describing the inter-neuronal relationships on a unit hypersphere. Our analysis reveals that maintaining this property is key to upholding the semantic generative capabilities of these models. To enhance stability during finetuning, we further introduce Constrained Orthogonal Finetuning (COFT), which supplements OFT with a constraint on the radius of the hypersphere. We focus our study on two prominent finetuning scenarios for text-to-image diffusion: subject-driven generation—where a few reference images and a prompt guide the synthesis of contextually varied images of the same subject—and controllable generation, where models are conditioned on supplementary signals. Experimental results demonstrate that our OFT approach achieves superior generation fidelity and faster convergence relative to competing techniques.
\end{abstract}

\section{Introduction}

Recent advances in text-to-image diffusion models~\cite{saharia2022photorealistic,ramesh2022hierarchical,rombach2022high} have significantly improved the realism and quality of images generated via textual guidance. Nonetheless, such guidance alone often lacks the precision required for nuanced and fine-grained control of generated content. To confront this limitation, we concentrate on two representative text-to-image generation tasks in this work:

\begin{itemize}[leftmargin=*,nosep]

    \item \textbf{Subject-driven generation}~\cite{ruiz2023dreambooth}: The objective is to synthesize new images of a specific subject placed in novel contexts, leveraging only a handful of subject images alongside a guiding text prompt.
    \item \textbf{Controllable generation}~\cite{zhang2023adding,mou2023t2i}: Given control information such as canny edges or segmentation masks, the aim is to direct the model to generate images adhering to this input in addition to a text prompt.
\end{itemize}

Both of these tasks ultimately focus on the challenge of finetuning text-to-image diffusion models in a way that maintains the generative capabilities acquired during pretraining. The desired properties for an effective finetuning approach can be summarized as follows: (1) \emph{training efficiency}, referring to a reduced number of trainable parameters and minimal training epochs; and (2) \emph{generalizability preservation}, ensuring that the high-quality and diverse outputs are retained. Common finetuning procedures involve either modifying existing neuron weights using a small learning rate (\emph{e.g.}, \cite{ruiz2023dreambooth}), or introducing a lightweight component with new parameterized neuron weights (\emph{e.g.}, \cite{hulora2022,zhang2023adding}). Although straightforward, neither method reliably safeguards the generative performance obtained during pretraining. Furthermore, there is a lack of systematic insight into the construction of effective finetuning schemes, as well as the selection of key hyperparameters such as training epoch count. One central challenge lies in the absence of a robust metric for assessing the extent to which pretrained generative abilities are retained. Current approaches frequently operate under the implicit belief that minimizing the Euclidean distance between the finetuned and pretrained models enhances the preservation of pretrained capabilities, which leads to the common use of very small learning rates. While this heuristic may sometimes apply, we contend that measuring only the Euclidean distance from the pretrained model does not sufficiently reflect the semantic fidelity retained post-finetuning. Therefore, introducing a more structurally meaningful metric to evaluate differences between the finetuned and pretrained models could significantly improve both the retention of pretraining performance and the robustness of the finetuning process.

Building on previous findings that hyperspherical similarity effectively represents semantic content~\cite{liu2017deep,liu2018decoupled,chen2020angular}, our approach leverages hyperspherical energy~\cite{liu2018learning} to describe the inter-neuronal relationships within a given layer. Hyperspherical energy, calculated as the aggregate of hyperspherical similarities (such as cosine similarity) across all neuron pairs within the same layer, reflects how uniformly neurons are distributed on the unit hypersphere~\cite{liu2021learning}. We posit that effective finetuning should yield only minimal changes in hyperspherical energy relative to the original pretrained model. A straightforward approach might entail imposing a regularization term that maintains hyperspherical energy throughout the finetuning process; however, this method does not guarantee an optimal minimization of energy differences. Instead, we exploit a crucial property: hyperspherical energy remains unchanged when a shared orthogonal transformation is applied to all neurons, due to invariance in pairwise hyperspherical similarity under such transformations. Leveraging this invariance, we introduce Orthogonal Finetuning~(OFT), a method for adapting extensive text-to-image diffusion models to target tasks without altering their hyperspherical energy. The primary mechanism of OFT involves learning an orthogonal transformation, applied uniformly across a layer, that preserves the pairwise angular relationships among neurons. This process can be interpreted as a reorientation of the underlying coordinate system for neurons within the layer. Additionally, acknowledging that reduced Euclidean distance between pretrained and finetuned models typically enhances retention of pretraining effectiveness, we present Constrained Orthogonal Finetuning~(COFT): an extension of OFT that restricts the finetuned neurons to reside on a hypersphere of fixed radius, centered at their original pretrained positions.

The rationale behind using orthogonal transformation in neuron finetuning draws partial inspiration from the 2D Fourier transform, a technique that represents an image in terms of its magnitude and phase spectra. Critically, the phase spectrum—which encodes the angular relationship between inputs and basis vectors—retains the bulk of semantic content in the image. Notably, even when combining the phase spectrum of a particular image with an arbitrary magnitude spectrum, the reconstructed image generally preserves its original semantics. This observation implies that modulating neuron orientations serves as an essential mechanism for imparting semantic modifications to generated images, which lies at the heart of both subject-driven and controllable generation tasks. Nonetheless, introducing unrestrained flexibility in adjusting neuron directions can disrupt the generative capabilities inherited from pretraining. To address this, we advocate for constraining these modifications by maintaining the pairwise angles between neurons. This approach is grounded in the hypothesis that these inter-neuronal angles play a pivotal role in encoding a neural network’s learned knowledge. Consequently, a logical step is to develop a layer-wise orthogonal transformation—shared within each layer—that preserves the hyperspherical energy, thereby leaving the angular relationships undisturbed.

Additionally, our methodology is informed by insights from orthogonal over-parameterized training~\cite{liu2021orthogonal}, wherein classification networks are constructed by applying orthogonal transformations to a randomly initialized neural architecture. The rationale stems from the observation that a network with random initialization typically possesses low expected hyperspherical energy. The objective in \cite{liu2021orthogonal} is to ensure this low-energy state is retained throughout training, given that reduced energy correlates with improved generalization in classification tasks~\cite{liu2018learning,lin2020regularizing}. They demonstrate that orthogonal transformations afford sufficient adaptability for training robust classifiers. Our work, in contrast, centers on finetuning text-to-image diffusion models to enhance controllability and boost generative performance for downstream tasks. There are two principal distinctions between OFT and \cite{liu2021orthogonal}. First, while the latter is designed to drive hyperspherical energy to a minimum, OFT instead endeavors to preserve the hyperspherical energy level characteristic of the pretrained model. This preservation prevents alteration of the essential semantic structures established during pretraining—a critical concern during the finetuning of diffusion models, as minimizing hyperspherical energy here could be detrimental. Second, OFT introduces constraints specifically to limit divergence from the pretrained parameters, resulting in a constrained version of the method; this stands in contrast to \cite{liu2021orthogonal}, where no such restriction is applied. Achieving an optimal balance between flexibility and stability is fundamental to successful finetuning, and we posit that OFT attains this balance effectively. Our main contributions can be summarized as follows:

\begin{itemize}[leftmargin=*,nosep]

    \item We introduce Orthogonal Finetuning, a new approach for refining text-to-image diffusion models that enhances controllability. To promote even greater training stability, we present a constrained version of this method which enforces a restriction on the angular shift relative to the original pretrained parameters.
    \item Unlike conventional finetuning techniques, OFT updates the model parameters while provably conserving the hyperspherical energy. Through empirical investigation, we demonstrate that this quantity is a crucial indicator of how well the semantic knowledge of the original generative model is retained.
    \item We evaluate OFT on two main tasks: subject-driven generation and controllable generation. Our thorough experimental analysis reveals substantial gains compared to existing approaches, particularly regarding output quality, speed of convergence, and stability during finetuning. Additionally, OFT displays superior sample efficiency, attaining satisfactory convergence with significantly fewer training images and epochs.
    \item To assess controllable generation, we propose a novel control consistency metric. The metric works by estimating the control signal from the output image and comparing it to the initial control signal. Our evaluations using this metric confirm the strong controllability achieved by OFT.
\end{itemize}

\section{Related Work}

\textbf{Text-to-image diffusion models}. Significant advancements~\cite{saharia2022photorealistic,ramesh2022hierarchical,rombach2022high,gu2022vector,nichol2022glide} have occurred in the field of text-to-image generation, primarily driven by breakthroughs in diffusion-based generative architectures~\cite{ho2020denoising,song2021score,song2021denoising,dhariwal2021diffusion} and joint vision-language representation learning~\cite{su2020vlbert,lu2019vilbert,radford2021learning,wang2021simvlm,li2022blip,li2023blip,alayrac2022flamingo,schuhmann2022laion}. For instance, GLIDE~\cite{nichol2022glide} and Imagen~\cite{saharia2022photorealistic} focus on diffusion modeling in pixel space. GLIDE performs joint optimization of the text encoder and diffusion prior using paired data, whereas Imagen relies on a pretrained, fixed text encoder. In contrast, latent-space approaches are used by Stable Diffusion~\cite{rombach2022high} and DALL-E2~\cite{ramesh2022hierarchical}. Here, Stable Diffusion employs VQ-GAN~\cite{esser2021taming} for constructing a visual codebook, while DALL-E2 utilizes CLIP~\cite{radford2021learning} to create a shared embedding space for both images and texts. Besides diffusion models, alternative generative frameworks such as generative adversarial networks~\cite{reed2016generative,zhang2017stackgan,xu2018attngan,li2019controllable} and autoregressive models~\cite{ramesh2021zero,ding2021cogview,wu2022nuwa,yuscaling} have also been explored for text-conditioned image synthesis. Notably, OFT is a model-agnostic strategy that is compatible with any text-to-image diffusion system.

\textbf{Subject-driven generation}. To avoid changes to the central subject, approaches such as~\cite{avrahami2022blended,nichol2022glide} incorporate user-provided masks as an auxiliary constraint. Inversion-based techniques~\cite{choi2021ilvr,dhariwal2021diffusion,ramesh2022hierarchical,gal2022image} enable context modifications while maintaining the integrity of the main subject. The method in~\cite{hertz2022prompt} allows for both localized and global editing even without the use of explicit masks. However, these strategies frequently fall short in preserving detailed identity information. Pivotal Tuning~\cite{roich2022pivotal} finetunes the generator near an inverted latent code, with an added regularization to maintain subject identity. In a related vein,~\cite{nitzan2022mystyle} learns a personalized face prior using a dataset of images from a specific individual. While~\cite{casanova2021instance} can synthesize a variety of instance representations, it may neglect instance-specific features. DreamBooth~\cite{ruiz2023dreambooth} tailors text-to-image diffusion models using a unique token, a handful of subject images, and a combination of reconstruction and class-specific prior preservation losses. Building on DreamBooth, OFT replaces conventional finetuning using a small learning rate with an orthogonal transformation-based optimization of the model parameters.

\textbf{Controllable generation}. Image-to-image translation is a prototypical case of controllable generation, with much of the earlier literature relying on conditional generative adversarial networks~\cite{isola2017image,zhu2017toward,wang2018high,park2019semantic,choi2018stargan}. Recent advances have extended diffusion models for this purpose as well~\cite{saharia2022palette,voynov2022sketch,wang2022pretraining}. Innovations such as ControlNet~\cite{zhang2023adding} enable precise modulation of a pretrained diffusion model by finetuning it in response to extra control signals, attaining strong performance in generating controllable outputs. Similarly, T2I-Adapter~\cite{mou2023t2i} introduces finetuning strategies for diffusion models to increase control over generated content. In alignment with the experimental settings outlined by~\cite{zhang2023adding,mou2023t2i}, we utilize OFT to finetune pretrained diffusion models. Our method achieves superior control capabilities while requiring less labeled data and involving fewer trainable parameters. Crucially, OFT imposes no additional computational costs at inference.

\textbf{Model finetuning}. The trend of finetuning large-scale pretrained models for downstream tasks has grown considerably in recent years~\cite{devlin2018bert,brown2020language,he2022masked}. Among finetuning strategies, various adaptation techniques (\emph{e.g.}, \cite{hulora2022,houlsby2019parameter,pfeiffer2020adapterfusion}) have been extensively explored, particularly in natural language processing. OFT is closely related to LoRA~\cite{hulora2022}, which is predicated on imposing a low-rank constraint on additive weight modifications during finetuning. However, OFT diverges by employing a multiplicative, orthogonal transformation that is shared across layers to adapt neuron weights, thereby preserving the angular relationships between neurons in a given layer and promoting enhanced stability.

\section{Orthogonal Finetuning}

\subsection{Why Does Orthogonal Transformation Make Sense?}

We first motivate the suitability of orthogonal transformation for finetuning text-to-image diffusion models. To do so, we break the problem into two parts: (1) the rationale behind adjusting neuron angles (\emph{i.e.}, modifying direction), and (2) the justification for employing orthogonal transformation in the process of finetuning these angles.

To address the first question, we are guided by empirical findings from \cite{liu2018decoupled,chen2020angular}, which indicate that differences in angular features serve as an effective metric for quantifying the semantic gap. SphereNet~\cite{liu2017deep} demonstrates that constraining all neuron activations to a unit hypersphere during neural network training preserves, and may even enhance, both model capacity and generalizability. This suggests that neuron directionality encapsulates the critical information present in the data. To further underscore the significance of neuron angles, we design a controlled experiment (refer to Figure~\ref{fig:toy_ae}) using a basic convolutional autoencoder trained on a collection of flower images. During the training phase, we generate feature maps using the conventional inner product (where $z$ represents an output of the convolutional kernel $\bm{w}$ applied to input $\bm{x}$ via a sliding window). In evaluation, we contrast three distinct methods for feature map extraction: (a) the inner product as during training, (b) extracting solely magnitude information, and (c) isolating angular information. Results depicted in Figure~\ref{fig:toy_ae} reveal that angular features enable nearly perfect reconstruction of the input, while magnitude conveys negligible information. Importantly, we refrain from employing cosine activation (c) during training; only the inner product is used to learn the model. These findings strongly suggest that the angles or directions of neurons are principally responsible for encoding the semantic content of the images. Consequently, to alter image semantics, adjusting neuron direction is likely to be a particularly effective approach.

Turning to the second question, the most direct strategy for adapting neuron direction is to apply a uniform rotation or reflection to all neurons within the same layer, which intrinsically results in an orthogonal transformation. Although it is possible to consider more elaborate angular transformations that adjust each neuron's direction independently, our experiments indicate that orthogonal transformations offer an optimal balance between adaptability and structural constraint. Furthermore, evidence from \cite{liu2021orthogonal} suggests that orthogonal transformations are adequately expressive for neural network learning. To substantiate this, we conduct an experiment focusing on the regularization properties imparted by enforcing orthogonality. Utilizing the setup of ControlNet~\cite{zhang2023adding}, we carry out a controllable generation task and present the outcomes in Figure~\ref{fig:ortho}. The results demonstrate that the standard OFT reliably enables precise control after completion of training (at epoch 20), whereas OFT without orthogonality fails to synthesize plausible images or achieve control. These findings reinforce the critical role of orthogonality constraints in OFT.

Traditionally, the finetuning process involves employing gradient descent with a low learning rate to adjust a model’s parameters (or certain layers within it). This conservative learning rate serves to limit changes from the original pretrained weights, such that standard finetuning implicitly tries to keep $\thickmuskip=2mu \medmuskip=2mu \| \bm{M} - \bm{M}^0\|$ small, where $\bm{M}$ designates the updated parameters after finetuning, and $\bm{M}^0$ represents the pretrained values. While this soft constraint is intuitive, it may nonetheless allow excessive flexibility for large-scale models. In response, LoRA enforces an additional low-rank condition on the weight modifications, specifically: $\thickmuskip=2mu \medmuskip=2mu \text{rank}(\bm{M} - \bm{M}^0)=r'$ with $r'$ being a small, fixed integer. By contrast, OFT constrains the similarity between pairs of neurons via: $\thickmuskip=2mu \medmuskip=2mu \| \text{HE}(\bm{M})-\text{HE}(\bm{M}^0) \|=0$, where $\text{HE}(\cdot)$ is the hyperspherical energy computed over the weight matrix. To demonstrate, take a fully connected layer represented as $\thickmuskip=2mu \medmuskip=2mu \bm{W}=\{\bm{w}_1,\cdots,\bm{w}_n\}\in\mathbb{R}^{d\times n}$, with each column vector $\bm{w}_i\in\mathbb{R}^{d}$ corresponding to the $i$-th neuron, and pretrained weights indicated by $\bm{W}^0$. The layer produces output $\thickmuskip=2mu \medmuskip=2mu \bm{z}\in\mathbb{R}^n$ according to $\thickmuskip=2mu \medmuskip=2mu \bm{z}=\bm{W}^\top\bm{x}$, where input $\thickmuskip=2mu \medmuskip=2mu \bm{x}\in\mathbb{R}^d$. The OFT method may thus be interpreted as minimizing the difference in hyperspherical energy between the finetuned and pretrained weights:

\begin{equation}\label{eq:oft_principle}
\footnotesize
\min_{\bm{W}} \left\| \text{HE}(\bm{W})-\text{HE}(\bm{W}^0)\right\|~~\Leftrightarrow~~\min_{\bm{W}} \bigg{\|} \sum_{i\neq j} \|\hat{\bm{w}}_i-\hat{\bm{w}}_j\|^{-1}-\sum_{i\neq j} \|\hat{\bm{w}}^0_i-\hat{\bm{w}}^0_j\|^{-1}\bigg{\|}
\end{equation}

Here, $\thickmuskip=2mu \medmuskip=2mu \hat{\bm{w}}_i:=\bm{w}_i/\|\bm{w}_i\|$ is the normalized vector of neuron $i$, and hyperspherical energy is defined for $\bm{W}$ by $\thickmuskip=2mu \medmuskip=2mu \text{HE}(\bm{W}):= \sum_{i\neq j} \|\hat{\bm{w}}_i-\hat{\bm{w}}_j\|^{-1}$. One can see that the objective in Eq.~\eqref{eq:oft_principle} achieves its minimum value, zero, whenever the weight matrices $\bm{W}$ and $\bm{W}^0$ differ only by a rotation or reflection, that is, if $\thickmuskip=2mu \medmuskip=2mu \bm{W}=\bm{R}\bm{W}^0$ for some orthogonal matrix $\thickmuskip=2mu \medmuskip=2mu \bm{R}\in\mathbb{R}^{d\times d}$ (where $\bm{R}$ has determinant $1$ for rotation or $-1$ for reflection). The central insight of OFT is thus to restrict finetuning to learning layer-wise, shared orthogonal transformations acting on the neurons of each layer. More precisely, OFT involves optimizing an

\begin{equation}\label{eq:oft_general}
    \footnotesize
    \bm{z}=\bm{W}^\top\bm{x}=(\bm{R}\cdot\bm{W}^0)^\top\bm{x},~~~\text{s.t.}~\bm{R}^\top\bm{R}=\bm{R}\bm{R}^\top=\bm{I}
\end{equation}

Here, $\bm{W}$ represents the weight matrix after applying OFT-finetuning, while $\bm{I}$ stands for the identity matrix. The operational details of OFT are depicted in Figure~\ref{fig:block}. In the same vein as LoRA's zero initialization approach, it is crucial to configure OFT so that fine-tuning precisely starts from the pretrained weights $\bm{W}^0$. This is accomplished by initially setting the orthogonal matrix $\bm{R}$ to the identity matrix, thereby aligning the initial state of the fine-tuned model with the original pretrained weights. To maintain the orthogonality constraint on $\bm{R}$ during training, various differentiable orthogonalization methods from \cite{liu2021orthogonal,lezcano2019cheap} can be leveraged. Further in the text, we address how such orthogonality can be preserved efficiently from a computational standpoint.

\subsection{Efficient Orthogonal Parameterization}\label{sect:eff_ortho}

Although differentiable, conventional orthogonalization procedures like the Gram-Schmidt process are typically impractical for large-scale settings due to their substantial computational burden~\cite{liu2021orthogonal}. To overcome this, we opt for the Cayley parameterization, which constructs the orthogonal matrix in a computationally economical manner. Specifically, the matrix is formed as $\thickmuskip=2mu \medmuskip=2mu \bm{R}=(\bm{I}+\bm{Q})(I-\bm{Q})^{-1}$, where the matrix $\bm{Q}$ is skew-symmetric, i.e., $\thickmuskip=2mu \medmuskip=2mu \bm{Q}=-\bm{Q}^\top$. This method does introduce a mild restriction: Cayley parameterization only yields orthogonal matrices of determinant $1$, meaning $\bm{R}$ belongs to the special orthogonal group. Nevertheless, this restriction does not empirically hinder performance. However, even when adopting the Cayley transform, representing a full-size orthogonal matrix $\bm{R}$ can be prohibitive in parameter cost when $d$ is large. To improve efficiency, we suggest utilizing a block-diagonal structure for $\bm{R}$, composed of $r$ independent blocks. This gives:

\begin{equation}
    \footnotesize
    \bm{R}=\text{diag}(\bm{R}_1,\bm{R}_2,\cdots,\bm{R}_r)=\begin{bmatrix}
    \bm{R}_1\in \text{O}(\frac{d}{r}) &  &\\
    &\ddots & \\
    & & \bm{R}_r\in \text{O}(\frac{d}{r})
    \end{bmatrix}\in \text{O}(d)
\end{equation}

Here, $\text{O}(d)$ refers to the orthogonal group of degree $d$, where both $\thickmuskip=2mu \medmuskip=2mu \bm{R}\in\mathbb{R}^{d\times d}$ and, for each $i$, $\thickmuskip=2mu \medmuskip=2mu \bm{R}_i\in\mathbb{R}^{d/r\times d/r}$ are orthogonal matrices. In the particular case when $\thickmuskip=2mu \medmuskip=2mu r=1$, the block-diagonal orthogonal structure reduces to a conventional, fully unconstrained orthogonal matrix. For a general orthogonal matrix of size $\thickmuskip=2mu \medmuskip=2mu d\times d$, the number of independent parameters totals $\thickmuskip=2mu \medmuskip=2mu d(d-1)/2$, which leads to $\mathcal{O}(d^2)$ complexity. When the orthogonal matrix is constructed in a block-diagonal manner with $r$ blocks, the total parameter count becomes $\thickmuskip=2mu \medmuskip=2mu d(d/r-1)/2$, corresponding to a complexity of $\thickmuskip=2mu \medmuskip=2mu \mathcal{O}(d^2/r)$. Further efficiency can be achieved by making all block matrices identical, i.e., setting $\thickmuskip=2mu \medmuskip=2mu\bm{R}_i=\bm{R}_j$ for all $i\neq j$, which lowers the parameter complexity to $\mathcal{O}(d^2/r^2)$. Importantly, these reductions in parameter count do not affect the orthogonality of $\bm{R}$, and thus, the hyperspherical energy is fully retained.

To compare OFT and LoRA with respect to parameter efficiency, consider LoRA’s trainable parameter count, given a low-rank setting $r'$, as $\thickmuskip=2mu \medmuskip=2mu r'(d+n)$. If we choose both $r$ and $r'$ proportional to the neuron dimension, such as $\thickmuskip=2mu \medmuskip=2mu r=r'=\alpha d$ for some constant $0<\alpha\leq 1$, LoRA’s parameter complexity increases to $\mathcal{O}(d^2+dn)$, while for OFT it reduces to $\mathcal{O}(d)$. To concretely demonstrate this difference, suppose a weight matrix of dimensions $\thickmuskip=2mu \medmuskip=2mu 128\times 128$; under this setup, LoRA entails $2,048$ learnable parameters when $\thickmuskip=2mu \medmuskip=2mu r'=8$, whereas OFT needs only $960$ trainable parameters at $\thickmuskip=2mu \medmuskip=2mu r=8$, assuming no block sharing.

\subsection{Constrained Orthogonal Finetuning}

A further restriction on the flexibility of OFT can be introduced by enforcing the fine-tuned model to stay within a limited neighborhood around the original pretrained model. Specifically, Constrained Orthogonal Finetuning (COFT) proposes the following forward computation:

\begin{equation}\label{eq:coft_general}
    \footnotesize
    \bm{z}=\bm{W}^\top\bm{x}=(\bm{R}\cdot\bm{W}^0)^\top\bm{x},~~~\text{s.t.}~\bm{R}^\top\bm{R}=\bm{R}\bm{R}^\top=\bm{I},~\norm{\bm{R}-\bm{I}}\leq \epsilon
\end{equation}

In this formulation, $\bm{R}$ is not only constrained to be orthogonal, but also required to lie within an $\epsilon$-ball of the identity matrix. While the orthogonality can be handled by the Cayley parameterization discussed in Section~\ref{sect:eff_ortho}, ensuring the additional $\epsilon$-neighborhood constraint within the Cayley framework is challenging. To analyze the Cayley transform more closely, one can employ the Neumann series expansion to approximate $\thickmuskip=2mu \medmuskip=2mu \bm{R}=(\bm{I}+\bm{Q})(I-\bm{Q})^{-1

series converges with respect to the operator norm). As a result, it is permissible to incorporate the constraint $\thickmuskip=2mu \medmuskip=2mu\|\bm{R}-\bm{I}\|\leq\epsilon$ within the Cayley transform itself; this translates equivalently to imposing $\thickmuskip=2mu \medmuskip=2mu \|\bm{Q}-\bm{0}\|\leq\epsilon'$, where $\bm{0}$ is the zero matrix and $\epsilon'$ represents a different tolerance parameter from $\epsilon$. This reformulated constraint on $\bm{Q}$ is amenable to enforcement via projected gradient descent methods. To guarantee that the orthogonal matrix $\bm{R}$ is initially the identity, we simply set $\bm{Q}$ to the zero matrix at initialization. The COFT framework can thus be interpreted as implementing two clear constraints: maintaining minimal hyperspherical energy variation and explicitly restricting deviations from the pretrained model. While most existing finetuning approaches implicitly include the second constraint, COFT formalizes it. Although OFT achieves strong results, our observations indicate that COFT improves finetuning stability, attributable to its explicit constraint on parameter deviation. Figure~\ref{fig:eps_coft} illustrates the influence of $\epsilon$ on COFT's empirical performance. It is apparent that $\epsilon$ modulates the degree of adaptation during finetuning: as $\epsilon$ increases, the outcomes approach those of OFT, whereas reducing $\epsilon$ causes the COFT-finetuned model to more closely track the original pretrained text-to-image diffusion model.

\subsection{Re-scaled Orthogonal Finetuning}

We introduce a straightforward modification to the standard OFT, wherein a learnable scaling coefficient is assigned to each neuron. This extension is inspired by the observation that scaling neurons leaves the hyperspherical energy unchanged (since normalization eliminates magnitude effects). The forward computation is as follows: $\thickmuskip=2mu \medmuskip=2mu\bm{z}=(\bm{R}\bm{W}^0\bm{D})^\top\bm{x}$\footnote{\textbf{Errata}: The NeurIPS camera ready manuscript contains a typographical error in the forward pass for re-scaled OFT, stated as $\thickmuskip=2mu \medmuskip=2mu\bm{z}=(\bm{D}\bm{R}\bm{W}^0)^\top\bm{x}$. The actual implementation aligns with the correct version, so the experimental results in Appendix~\ref{app:rescaled} are unaffected.} where $\thickmuskip=2mu \medmuskip=2mu \bm{D}=\text{diag}(s_1,\cdots,s_n)\in\mathbb{R}^{n\times n}$ is a diagonal matrix of trainable parameters, with all $s_1,\ldots,s_n$ strictly positive. This approach differs from the original OFT formulation (Eq.~\eqref{eq:oft_general}), in which only $\bm{R}$ is optimized; here, both $\bm{R}$ and the scaling matrix $\bm{D}$ are trainable. The proposed re-scaled OFT variant enhances the modeling flexibility of OFT with only a minor increase in parameter count. For experimental clarity and to isolate the effects of the orthogonal transformation, we adopt the basic OFT for our main results, though empirical evidence suggests that re-scaled OFT frequently yields superior performance (see Appendix~\ref{app:rescaled}).

\section{Intriguing Insights and Discussions}

\textbf{OFT is architecture-independent}. The OFT approach can be theoretically employed with any neural network architecture. In the case of Transformers, LoRA is most commonly integrated into attention weight matrices~\cite{hulora2022}. To ensure a fair comparison with LoRA, we restrict the application of OFT in our experiments solely to attention weight matrices. In addition to dense (fully connected) layers, OFT is particularly compatible with convolutional layers, as the block-diagonal form of $\bm{R}$ has meaningful analogues in this context—unlike LoRA. Assigning each block to a different input channel means that each block performs a channel-specific transformation, reminiscent of learning depth-wise convolution filters~\cite{chollet2017xception}. If, instead, a single block is shared across all channels, the orthogonal transformation is uniformly applied across neurons in each channel. Additional exploratory experiments regarding convolution layer finetuning using OFT are detailed in Appendix~\ref{app:convolution}.

\textbf{Connection to LoRA}. LoRA achieves parameter efficiency by augmenting the pretrained weight matrix with a low-rank update, thereby preserving much of the original information. Conversely, OFT maintains the integrity of the pretrained weights by constraining the transformation matrix to be orthogonal (i.e., of full rank), thereby ensuring that the transformation does not overwrite learned information from pretraining. The forward computation in OFT can be reformulated as $\thickmuskip=2mu \medmuskip=2mu \bm{z}=(\bm{R}\bm{W}^0)^\top\bm{x}=(\bm{W}^0+(\bm{R}-\bm{I})\bm{W}^0)^\top\bm{x}$, where $\thickmuskip=2mu \medmuskip=2mu (\bm{R}-\bm{I})\bm{W}^0$ serves a comparable function to LoRA’s low-rank adaptation term. Since $\bm{W}^0$ generally possesses full rank, OFT can similarly effectuate a low-rank adjustment if $\thickmuskip=2mu \medmuskip=2mu \bm{R}-\bm{I}$ is chosen to be low-rank. In LoRA, the parameter $r'$ regulates the rank for efficiency; analogously, OFT uses the size of diagonal blocks, $r$, to reduce the trainable parameter count. Notably, LoRA and OFT exemplify two orthogonal strategies for efficient parameterization: LoRA utilizes low-rank structure for compactness, whereas OFT leverages block-diagonal sparsity with orthogonal constraints.

\textbf{Why OFT converges faster}. As indicated in Figure~\ref{fig:toy_ae}, adjusting the directionality of neurons proves most effective for altering model semantics, an operation which OFT is tailored to perform. Furthermore, OFT’s framework can be interpreted as neuron finetuning on a smooth hypersphere manifold, facilitating more favorable optimization conditions. This perspective is also supported empirically in \cite{liu2021orthogonal}.

\textbf{Why not minimize hyperspherical energy}. Unlike \cite{liu2021orthogonal}, our approach does not target minimizing hyperspherical energy. For classification tasks, non-redundant neuron representations are advantageous, and attaining minimum hyperspherical energy corresponds to an even dispersal of neuron directions on the hypersphere, which may actually undermine the pretrained structure. Therefore, minimizing this energy does not constitute a meaningful finetuning goal, as it can compromise essential pretrained knowledge.

\textbf{Balancing flexibility and regularity in finetuning.} Our findings indicate a fundamental compromise between model adaptability and stability. While conventional finetuning achieves maximal adaptability, it tends to lack robustness and is prone to destabilization or model degradation. OFT, despite its conceptual simplicity, achieves an effective equilibrium by maintaining the angular relationships between neuron vectors, thus imparting both adaptability and control. Moreover, the block-diagonal approach can be interpreted as imposing an even more pronounced regularization on the orthogonal transformation.

\textbf{Zero inference latency increase.} In contrast to methods such as ControlNet, our OFT approach does not add any computational burden during model inference. During deployment, the orthogonal matrix $\bm{R}$ learned during training can be directly incorporated into the original weights $\bm{W}^0$ by simple matrix multiplication, resulting in a new weight matrix $\thickmuskip=2mu \medmuskip=2mu \bm{W}=\bm{R}\bm{W}^0$, functionally indistinguishable in inference speed from the pretrained model.

\section{Experiments and Results}

\textbf{Experimental configuration.} We base our evaluations on Stable Diffusion v1.5~\cite{rombach2022high} for the pretrained text-to-image framework. To maintain impartiality, samples for all methods are drawn at random. Our protocol for subject-driven image synthesis aligns closely with DreamBooth~\cite{ruiz2023dreambooth}, while for conditional image generation, we adopt the approaches from ControlNet~\cite{zhang2023adding} and T2I-Adapter~\cite{mou2023t2i}. For consistent benchmarking against LoRA, we restrict the application of OFT or COFT to those specific layers employed by LoRA. Supplementary results and implementation specifics are elaborated in Appendix~\ref{app:settings}.

\subsection{Subject-driven Generation}

\textbf{Configuration details.} Baselines used in our comparisons include DreamBooth~\cite{ruiz2023dreambooth} and LoRA~\cite{hulora2022}. All compared models are trained with the same objective function as specified in DreamBooth. For both DreamBooth and LoRA, we adhere to protocols and hyperparameter choices as outlined in the original literature. Comprehensive results and additional insights are provided in Appendix~\ref{app:settings},\ref{app:coft_oft},\ref{app:more_qualitative},\ref{app:failure}.

\textbf{Finetuning stability and convergence}. To begin, we assess the robustness and rate of convergence during finetuning for DreamBooth, LoRA, OFT, and COFT, with outcomes illustrated in Figure~\ref{fig:teaser} and Figure~\ref{fig:db_convergence}. It is evident that both COFT and OFT facilitate stable finetuning of the diffusion model. After 400 steps, effective control is achieved with both DreamBooth and OFT approaches, whereas LoRA does not succeed in maintaining subject identity. Following 2000 steps, DreamBooth displays image collapse, and LoRA generates images with yellow fur rather than the intended yellow shirt. Conversely, OFT and COFT continue to provide consistent and reliable manipulation of the generated outputs. These observations support the assertion that our OFT approach delivers rapid yet reliable convergence for subject-driven generative tasks. Importantly, the methods' robustness to the number of finetuning iterations is a significant advantage, as it reduces the need for meticulous hyperparameter selection for different subjects. OFT and COFT allow for the straightforward assignment of a large iteration count without extensive tuning. In particular, with COFT and an appropriate $\epsilon$, both the learning rate and number of iterations become easy to configure.

\textbf{Quantitative comparison}. In line with \cite{ruiz2023dreambooth}, we provide a quantitative assessment measuring subject preservation (DINO~\cite{caron2021emerging}, CLIP-I~\cite{radford2021learning}), text prompt alignment (CLIP-T~\cite{radford2021learning}), and diversity among samples (LPIPS~\cite{zhang2018unreasonable}). CLIP-I calculates the mean pairwise cosine similarity between CLIP feature representations of generated and real images. DINO adopts a similar process, but leverages ViT S/16 DINO embeddings. CLIP-T refers to the average cosine similarity between text prompt embeddings and those of generated images via CLIP. Additionally, we report the average LPIPS cosine similarity between images generated for the same subject and prompt. As shown in Table~\ref{table:dreambooth_final}, both COFT and OFT significantly surpass DreamBooth and LoRA in terms of DINO and CLIP-I scores, and display marginally superior or competitive results regarding prompt adherence and sample diversity. For every method evaluated, we perform finetuning of the identical pretrained model across 30 distinct random seeds to mitigate variance due to randomness. The outcomes demonstrate that OFT not only exhibits enhanced convergence and stability but also consistently achieves higher-quality results.

\textbf{Qualitative comparison}. To provide a more direct illustration of OFT's advantages, we display a selection of randomly chosen outputs for subject-driven generation in Figure~\ref{fig:dreambooth_final}. All samples are generated using the same finetuned base model per method to ensure a fair evaluation; in particular, no individual prompt engineering is conducted for optimal image results. For each technique, the model checkpoint that attained the top validation CLIP metrics is presented. As observed in Figure~\ref{fig:dreambooth_final}, OFT and COFT consistently retain the semantic attributes of the subject, whereas LoRA frequently fails to preserve subject identity (for instance, the LoRA model loses all recognizable features of the subject in the bowl example). Additionally, both OFT and COFT demonstrate superior precision in following the text prompts, whereas DreamBooth, while adept at subject identity retention, often struggles to faithfully generate content as described by the prompt (as seen in the initial row of the bowl example). This qualitative assessment underscores that the OFT approach strikes an effective balance between text-driven controllability and preservation of subject identity. Furthermore, OFT’s performance is largely robust to the number of training iterations, maintaining strong results even at higher iteration counts, in contrast to DreamBooth and LoRA. Additional qualitative examples can be found in Appendix~\ref{app:more_qualitative}. We also report a human evaluation in Appendix~\ref{app:human}, which further substantiates the advantages of OFT.

\subsection{Controllable Generation}

\textbf{Settings}. As comparative baselines, we employ ControlNet~\cite{zhang2023adding}, T2I-Adapter~\cite{mou2023t2i}, and LoRA~\cite{hulora2022}. The main text addresses three difficult tasks for controllable image generation: Canny edge to image~(C2I) utilizing the COCO dataset~\cite{lin2014microsoft}, segmentation map to image~(S2I) on ADE20K~\cite{zhou2017scene}, and landmark to face~(L2F) leveraging CelebA-HQ~\cite{karras2018progressive,xia2021tedigan}. For all methods, Stable Diffusion~(SD) v1.5 is finetuned on the respective datasets for 20 epochs. Supplementary results are provided in Appendix~\ref{app:more_qualitative}, \ref{app:more_control_task}, and \ref{app:failure}.

\textbf{Convergence}. To assess the convergence characteristics of ControlNet, T2I-Adapter, LoRA, and COFT in the context of the L2F task, we conduct both quantitative and qualitative analyses. For quantitative assessment, we measure the average $\ell_2$ distance between the control face landmarks and those predicted by each method. Figure~\ref{fig:face_convergence} illustrates the evolution of face landmark error as a function of the number of finetuning epochs. The results demonstrate that COFT and OFT converge markedly faster than the other approaches. Notably, LoRA requires 20 epochs to reach the level of performance that our OFT framework achieves by the 8-th epoch. Importantly, the trainable parameter count in OFT and COFT closely matches that of LoRA (and is considerably lower than that of ControlNet), yet our methods converge with greater efficiency than the current state-of-the-art. Furthermore, the rapid convergence of OFT is substantiated by the observations in Figure~\ref{fig:teaser}, where the example on the right highlights OFT’s superior data efficiency: it achieves satisfactory convergence using just 5\% of the ADE20K dataset—significantly less data than required by ControlNet or LoRA. 

Turning to qualitative evaluation, we emphasize the comparison among OFT, COFT, and ControlNet, as ControlNet produces landmark errors most comparable to those of our methods. Figure~\ref{fig:face_convergence_qual} presents that both OFT and COFT achieve stable convergence, progressively producing face poses that align well with the specified control landmarks. Unlike our frameworks, where convergence occurs smoothly and gradually, ControlNet’s controllability tends to emerge abruptly after the 8-th epoch. This observation is consistent with the sudden convergence behavior previously reported in \cite{zhang2023adding}. The inherent stability and smoothness of OFT’s convergence not only facilitate practical deployment but also render the process of hyperparameter selection more straightforward, owing to the more predictable and manageable training dynamics.

\textbf{Quantitative comparison}. To assess the effectiveness of controllable generation, we propose a metric that quantifies control consistency. This metric involves extracting the control signal from the output image and comparing it against the original input control signal. For the C2I task, we utilize IoU and F1 score as evaluation measures. In the context of S2I, we report mean IoU, as well as mean and overall accuracy. For L2F, the evaluation is based on the mean $\ell_2$ distance between predicted and control landmarks. Additional specifics on these consistency metrics can be found in Appendix~\ref{app:settings}. For each method under comparison, optimal hyperparameter configurations are employed. The data in Table~\ref{table:control_gen} indicates that both OFT and COFT achieve superior and more precise control relative to other baselines. We note that methods based on adapters (\emph{e.g.}, T2I-Adapter and ControlNet) are slower to converge and exhibit inferior performance at convergence. Although LoRA outperforms ControlNet in the S2I scenario, it falls short in C2I and L2F tasks. Broadly, our experiments suggest that LoRA’s performance is capped at a relatively modest level, even after extensive hyperparameter optimization. In contrast, the OFT framework continues to improve as the number of trainable parameters increases, showing no evidence of plateauing. We highlight that our quantitative evaluation approach for controllable generation is a novel aspect of this work, providing a reliable measure of control efficacy for finetuned models (contingent on the accuracy of the used segmentation/detection models).

\textbf{Qualitative comparison}. We conduct a qualitative assessment of OFT and COFT in relation to several leading approaches such as ControlNet, T2I-Adapter, and LoRA. As illustrated in the randomly sampled outputs in Figure~\ref{fig:control_final}, both OFT and COFT consistently deliver not only high-quality and photorealistic results but also precise controllability over the generated images. In the S2I scenario, LoRA is unable to faithfully replicate the target segmentation map, whereas ControlNet, OFT, and COFT manage to closely adhere to the input structure. Notably, compared to ControlNet, OFT and COFT succeed in synthesizing visually compelling images featuring enhanced detail and more coherent geometric composition, despite utilizing substantially fewer parameters. Regarding the C2I task, OFT and COFT effectively generate images that align semantically with coarse Canny edges, outperforming both T2I-Adapter and LoRA. For the L2F task, our approach achieves superior accuracy in pose guidance, particularly for faces with challenging orientations. Across all three control scenarios, OFT and COFT consistently outperform existing methods, substantiating the superiority of the OFT framework for controllable image generation. Further extensive qualitative results for each control setting are available in Appendix~\ref{app:control_exp}; Appendix~\ref{app:more_control_task} also demonstrates OFT’s generalizability to additional control modalities, including dense pose-to-human body, sketch-to-image, and depth-to-image translations.

\section{Concluding Remarks and Open Problems}

Prompted by the recognition that the angular relationships among neurons play a pivotal role in shaping visual semantics, we introduce a straightforward yet powerful finetuning approach—orthogonal finetuning—to enable greater control over text-to-image diffusion models. Our approach is applied specifically to two use cases: subject-driven generation and controllable generation. In comparison to prevailing techniques, OFT achieves enhanced controllability and finetuning stability, utilizing a smaller set of tunable parameters. A further significant advantage is that OFT imposes no extra computational burden during inference, ensuring that models remain practical for real-world deployment.

OFT gives rise to several notable unresolved questions. Firstly, the orthogonality constraint in OFT is maintained via Cayley parametrization, which entails a matrix inversion step that can be a bottleneck for scalability. While our adoption of block diagonal parametrization partially mitigates this issue, developing differentiable strategies for accelerating this inversion is an ongoing challenge. Secondly, OFT holds particular promise for compositionality: orthogonal matrices derived from multiple OFT finetuning processes can be composed through multiplication and still maintain orthogonality. Investigating whether this collection of orthogonal matrices continues to encode knowledge from all individual downstream tasks offers a compelling avenue for future research. Lastly, the method's parameter efficiency is closely tied to the block diagonal design, a choice that may impose certain biases and restrict adaptability. Discovering ways to boost parameter efficiency with fewer biases and greater effectiveness stands as a key open problem in this area.

\end{document}